\subsection{Diagnostic levers}

Each lever below attaches to exactly one factor of the decomposition introduced earlier (annotated in the Factor column); the table is followed by three refutable propositions that localise \emph{why}, not merely \emph{where}, recall is lost.

\begin{table}[h]
\centering
\footnotesize
\begin{tabular}{p{2.6cm}p{2.3cm}p{6.3cm}p{2.6cm}}
\toprule
Lever & Factor & Finding & Evidence \\
\midrule
\textbf{Learned vs.\ prior} \emph{(\texttt{top\_k}-limited rule-selection probe --- not deployed recall)} &
\texttt{selection\_\allowbreak retention} &
A SyGMa-style frequency prior significantly out-ranks the learned generator: gen-only recall@15 prior-only \textbf{0.410} vs learned-only \textbf{0.266} ($\Delta$ \textbf{$-$0.144}, 95\% CI \textbf{[$-$0.196, $-$0.095]}, paired bootstrap, n=245); with the filter, 0.405 vs 0.300 ($\Delta$ $-$0.105, CI [$-$0.152, $-$0.058]). Adding the prior to the learned scorer lifts \textbf{+0.130} gen / \textbf{+0.099} filter (both significant); the filter significantly helps only the \emph{weak learned} ordering (+0.034, CI [+0.011, +0.061]), not the strong prior (n.s.). &
\texttt{results/\allowbreak prior\_\allowbreak vs\_\allowbreak learned.json} (245 test subs) \\
\addlinespace
\textbf{Multi-step (depth-2)} &
\texttt{coverage\allowbreak\_bank} &
Breadth-capped depth-2 rule application lifts the ceiling by only \textbf{+0.012} (0.711$\to$0.723) at \textbf{8.5$\times$} candidate cost (194$\to$1653 candidates/substrate) --- not the dominant coverage lever. &
\texttt{results/\allowbreak benchmark\_\allowbreak report\_\allowbreak depth2.json} (150 subs, beam 10) \\
\addlinespace
\textbf{Coverage ($\Delta$MW gap)} &
\texttt{coverage\allowbreak\_bank} &
\textbf{26.9\%} of true metabolites are uncovered by the depth-1 bank (plain InChIKey; $\sim$25\% under tautomer). Misses are a diverse long tail --- the top missing class (hydroxylation/oxidation) is only \textbf{6\%} of uncovered. &
\texttt{results/\allowbreak benchmark\_\allowbreak report\_\allowbreak gap.json} (500 subs) \\
\addlinespace
\textbf{Data scaling} &
\texttt{selection\_\allowbreak retention} &
Recall saturates (2418$\to$4787 substrates $\approx$ flat) --- the plateau is not a data-quantity problem. &
\texttt{results/\allowbreak full\{2500,\allowbreak 5000\}\allowbreak \_single.log} \\
\addlinespace
\textbf{Regioselectivity (SoM)} &
\texttt{ranking\allowbreak\_conversion} &
A site-of-metabolism prior gives only a small lift --- regioselective ranking is hard within the bank. &
\texttt{results/\allowbreak train\_\allowbreak som.log} \\
\addlinespace
\textbf{Selection breadth (\texttt{top\_k} sweep)} &
\texttt{selection\_\allowbreak retention} &
Widening the deployed rule selector's \texttt{top\_k} 30$\to$300 lifts recall@15 \textbf{0.352$\to$0.413 (+0.061)} and pool coverage \textbf{0.482$\to$0.608}, saturating below the 0.735 ceiling (n=245 test subset; trend, not deployed-scale absolute). &
\texttt{results/\allowbreak selection\_\allowbreak ablation.json} (245 test subs) \\
\addlinespace
\textbf{Abstention (filter gate)} &
\texttt{ranking\allowbreak\_conversion} &
Gating output by a filter-score threshold $\tau$ traces a recall$\leftrightarrow$precision frontier, but precision stays \textbf{nearly flat ($\sim$0.10--0.12)} and no val-selected $\tau$ reaches 0.2: raising $\tau$ shrinks the output (9.96$\to$0.3) and collapses recall (\textbf{0.388$\to$0.041}) for essentially no precision gain --- precision is annotation-bounded, not threshold-bounded, justifying the rank-only default. &
\texttt{results/\allowbreak abstention\_\allowbreak frontier.json} (val-selected $\to$ 1170 test) \\
\bottomrule
\end{tabular}
\caption{Diagnostic levers and the recall-decomposition factor each one attaches to.}
\label{tab:levers}
\end{table}

\subsection{Selection breadth and abstention ablations}

\paragraph{Selection-breadth counterfactual (n=245 test subset).}
Sweeping the number of rules applied (\texttt{top\_k}) on the deployed pipeline (generator from \texttt{full5000\_priors} with the trained prior restored, deployed filter, ranked by the deployed \texttt{filter} $\times$ \texttt{generator} score) isolates rule-selection breadth as its own lever:

\begin{table}[h]
\centering
\begin{tabular}{lcccc}
\toprule
\texttt{top\_k} (rules applied) & pool coverage (oracle) & deployed recall@15 & mean pool & mean output \\
\midrule
30 (deployed) & 0.482 & 0.352 & 12.5 & 7.99 \\
100 & 0.547 & 0.388 & 35.3 & 9.47 \\
300 ($\approx$ applicable-rule limit) & 0.608 & 0.413 & 107.6 & 9.53 \\
\bottomrule
\end{tabular}
\caption{Selection-breadth counterfactual on the deployed pipeline (n=245 test subset).}
\label{tab:selection-breadth}
\end{table}

Widening \texttt{top\_k} 30$\to$300 monotonically lifts recall@15 \textbf{0.352$\to$0.413 (+0.061)} and pool coverage \textbf{0.482$\to$0.608}, at the cost of an \textbf{8.6$\times$} pool inflation (12.5$\to$107.6 mean candidates) --- recall bought with precision, the same trade SyGMa makes by overproducing ($\sim$74 outputs/substrate). It also saturates below the ceiling: pool coverage plateaus at \textbf{0.608 $<$ 0.735} even at $\sim$all applicable rules, and recall@15 0.413 still $\ll$ SyGMa's 0.572 --- so breadth is a real but partial lever. Ranking still loses roughly a third of in-pool hits even at \texttt{top\_k}=300 (recall / pool coverage = 0.413 / 0.608 = \textbf{0.68}, consistent with \texttt{ranking\allowbreak\_conversion} = 0.726, reported in the decomposition-results section).

\paragraph{Abstention is not a precision lever (val-selected $\tau$, full 1170-substrate test).}
Where breadth trades precision for recall, the reverse move --- deliberately abstaining --- barely moves precision. Gating the deployed pipeline's output by a filter-score threshold $\tau$ (candidates with \texttt{filter\_score} $<$ $\tau$ are dropped; survivors ranked by \texttt{filter} $\times$ \texttt{gen}, top-15) and sweeping $\tau$ traces a recall$\leftrightarrow$precision frontier whose precision axis is \textbf{nearly flat}: from rank-only $\tau=0$ (recall@15 \textbf{0.388}, precision \textbf{0.098}, 9.96 outputs/substrate) through $\tau=0.5$ (recall 0.146, precision 0.114, 1.55 outputs) to $\tau=0.9$ (recall 0.041, precision 0.040, 0.30 outputs), precision never exceeds \textbf{$\sim$0.123}, and no $\tau$ selected on validation reaches even precision 0.2. Raising $\tau$ removes true and apparently-false candidates in near-equal proportion --- because most ``false'' predictions are \emph{unannotated true} metabolites --- so gating shrinks the output and collapses recall for almost no precision gain; the best validation-selected F1 operating point ($\tau \approx 0.10$) lifts F1 only marginally (0.129$\to$0.131) by trimming the output tail. Precision here is therefore \textbf{annotation-bounded, not threshold-bounded}, which both substantiates the ``precision is a pessimistic lower bound under incomplete annotation'' caveat (noted in the evaluation-hygiene and decomposition-results sections) and justifies the rank-only deployment default (\texttt{results/\allowbreak abstention\_\allowbreak frontier.json}).

On the same broad pools, a second arm compares the deployed learned ranker (\texttt{filter} $\times$ \texttt{gen}) against the rule frequency prior alone (SyGMa-style):

\begin{table}[h]
\centering
\begin{tabular}{lcc}
\toprule
\texttt{top\_k} & filter $\times$ gen (learned) & frequency prior (SyGMa-style) \\
\midrule
30 & 0.352 & 0.359 \\
100 & 0.388 & 0.374 \\
300 & 0.413 & 0.374 \\
\bottomrule
\end{tabular}
\caption{Learned ranker vs.\ frequency prior on identical broad pools, by selection breadth.}
\label{tab:ranker-vs-prior}
\end{table}

At narrow breadth the two are $\approx$ equal (0.352 vs 0.359); as the pool broadens the \textbf{learned filter pulls ahead} (0.413 vs 0.374 at \texttt{top\_k}=300), while the frequency prior plateaus at 0.374 --- swapping to the prior signal at breadth would hurt, not help.

\emph{Synthesis:} GRAIL's deficit vs.\ SyGMa on this axis is a compound of (i) over-aggressive rule selection at the deployed \texttt{top\_k}=30 (recoverable $\sim$+0.06 by widening breadth) and (ii) a residual rule-set coverage gap (pool coverage saturates at 0.608 $<$ 0.735) --- it is \textbf{not} a ranking-model deficiency: the learned filter out-ranks the frequency prior once the candidate pool is broad. This is an \textbf{n=245 test-subset} ablation; the trend (breadth helps, ranking doesn't hurt) is the finding --- the absolute \texttt{top\_k}=30 baseline (0.352) is subset-inflated relative to the full-1170 deployed macro recall (0.330). (\texttt{results/\allowbreak selection\_\allowbreak ablation.json}, \texttt{results/\allowbreak selection\_\allowbreak ablation\_\allowbreak ranksignal.json})

\paragraph{Probe vs.\ deployed.}
The learned-vs-prior row is a controlled, \texttt{top\_k}-limited rule-selection \textbf{probe} (each mode picks its own top-30 rules, same downstream product loop and filter) --- it is \emph{not} the deployed recall. The deployed pipeline in fact applies its \textbf{top-30 selected rules} --- the selection-breadth ablation above confirms this: its \texttt{top\_k}=30 baseline (0.352 subset) reproduces deployed recall (0.330 macro / 0.344 3-seed, reported in the decomposition- and deployment-results sections), so the deployed pipeline \emph{is} a \texttt{top\_k}=30 rule selector, and widening that selector's breadth is itself a soft lever (previous paragraph), not evidence that rules are applied unselectively. Separately --- and narrower in scope --- the deployed pipeline is measurably \textbf{prior-independent} at that fixed \texttt{top\_k}=30 operating point: on an earlier 291-substrate evaluation, two checkpoints differing only in whether a trained prior buffer was present gave essentially identical recall (0.335 vs 0.334), so deployed recall (\textbf{0.330} macro / \textbf{0.261} micro, reported in the decomposition- and deployment-results sections) is coverage-limited, not prior-limited at the deployed operating point --- the prior \emph{term} barely moves recall there, distinct from the \texttt{top\_k} breadth axis above; the probe \emph{reveals} the learned selector's weakness; deployment \emph{masks} it by not depending on that selector's ranking. \emph{Honesty note:} an earlier draft reported the opposite --- ``learned beats prior'' --- an artifact of a checkpoint whose prior buffer was un-persisted; that reversal is withdrawn and corrected here, caught by adversarial verification. Neither multi-step application nor any single rule-family addition moves coverage much (misses are a diverse long tail), and data scaling is flat --- so the dominant, addressable loss remains \texttt{selection\_\allowbreak retention} (decomposition-results section: 0.489), the learned rule-selector performing worse than a trivial frequency prior in the probe. The three propositions below explain \emph{why}.

\subsection{Three propositions}

\paragraph{Proposition 1 --- Surrogate mismatch ($\to$ \texttt{ranking\allowbreak\_conversion}).}
A filter trained by a strictly proper scoring rule (BCE/PU) learns a globally calibrated posterior, Bayes-optimal for AUC and calibration; ranking each substrate's candidate pool by that posterior and taking top-$k$ is recall@$k$-optimal only when pools are homogeneous. GRAIL's pools vary in size (17--150) and positive rate (\texttt{n\_true} 1--18), so a pointwise-calibrated scorer can be recall@$k$-suboptimal even while a listwise, ranking-consistent surrogate dominates --- supported by a minimal 2-substrate counterexample (\texttt{grail\_\allowbreak metabolism/\allowbreak tests/\allowbreak test\_\allowbreak prop1\_\allowbreak counterexample.py}) in which the recall-superior reorder is verified \emph{not} globally calibrated, so a proper-scoring objective rejects it.

\emph{Confirmation --- a controlled objective swap on the full test set.} The cleanest evidence isolates the objective while holding everything else fixed: the same two auxiliary heads, the same frozen generator and filter, and the \textbf{same candidate pool} are used twice, differing only in how the heads were trained. Trained \textbf{pointwise} (independent MLE) and multiplied into the rank, they add \textbf{nothing} --- +0.0002 (95\% CI [$-$0.0073, +0.0077], n.s.). Fine-tuned instead with a \textbf{listwise ranking loss} against the frozen $\log(\text{filter})+\log(\text{gen})$ context, the identical heads gain \textbf{+0.0089 (95\% CI [+0.0032, +0.0153])} over their pointwise counterpart and \textbf{+0.0091 (95\% CI [+0.0027, +0.0160])} over the \texttt{filter} $\times$ \texttt{gen} baseline --- paired bootstrap over the full clean test, n=1170 (\texttt{results/\allowbreak joint\_\allowbreak rerank.json}). Since capacity, features, and pool are held constant, the margin is attributable to the \textbf{training objective}, which is exactly the surrogate mismatch this proposition predicts. A second, independent re-ranking arm reproduces the direction at the same n: \texttt{filter} $\times$ \texttt{gen} $\times$ \texttt{type} $\times$ \texttt{site} lifts recall@15 from 0.388 to 0.404, a paired \textbf{+0.0165 (95\% CI [+0.006, +0.027])}, while the factorized signal on its own is n.s. (+0.001, [$-$0.010, +0.013]) (\texttt{results/\allowbreak hybrid\_\allowbreak rerank\_\allowbreak full1170.json}).

\emph{Secondary, weaker evidence:} a higher-capacity listwise-InfoNCE reranker scored \textbf{0.433 $\to$ 0.500 @15 (+0.067)}, 74\% of the Stage-2 oracle 0.677, on a held-out Stage-2 run (\texttt{docs/\allowbreak benchmark/\allowbreak stage2\_\allowbreak ranker\_\allowbreak evidence.md}, Spike-3) --- reported as $\pm$0.015 across 3 seeds (pure reranker variance; the generator baseline is deterministic). We do \textbf{not} give this figure a paired interval: its per-substrate scores were not retained, so a paired bootstrap cannot be recomputed without retraining, and the proposition rests on the n=1170 intervals above rather than on this number. \emph{Guardrail:} this is a theorem about \textbf{objectives}, not a recall win --- the reranker's \textbf{0.500 still loses to SyGMa (0.558)} and is reported only as a separate Stage-2 artifact, never as a headline number.

\paragraph{Proposition 2 --- Learning loses to counting: rule selection is extreme multi-label classification ($\to$ \texttt{selection\_\allowbreak retention}).}
This is the paper's central selection finding, and it is not metabolism-specific folklore rediscovered: selecting which of the \textbf{7,581} rules apply to a substrate is an \emph{extreme multi-label classification} (XMC) problem with a near-empty positive set ($\approx$1--3 firing rules per substrate, $\sim$0.03\% of labels). In exactly this regime, XMC has long established that frequency/propensity baselines are hard to beat and that na\"{\i}ve losses are dominated by label sparsity --- propensity-scored losses were introduced precisely for this failure mode \citep{Jain_2016}. Metabolism rule selection \textbf{inherits this XMC pathology}, which the domain literature does not name; our contribution is to identify it here and to test it causally (below). Concretely: under PU annotation with an approximately constant labeling propensity (SCAR), a learner with constant unlabeled weighting recovers a propensity-distorted score whose dominant component is the marginal rule-firing rate $\pi(r)$ --- the frequency prior --- so the prior is Bayes-competitive \textbf{by construction} at the rule-selection stage (which rules clear \texttt{top\_k}), and the learned selector improves on it only if substrate-conditional prevalence variation exceeds estimation noise; the observed data-scaling saturation (Table~\ref{tab:levers}, row 4) indicates it does not at current scale.

\emph{Anchor:} learned-only \textbf{0.266} vs prior-only \textbf{0.410} (gen-only @15, $\Delta$ \textbf{$-$0.144}, 95\% CI \textbf{[$-$0.196, $-$0.095]}, paired bootstrap, n=245; \texttt{results/\allowbreak prior\_\allowbreak vs\_\allowbreak learned.json}). \emph{Ranking-stage nuance:} this prior advantage is selection-specific, not a general ``learned $<$ prior'' --- at the separate \textbf{ranking} stage (ordering the candidate pool once rules are already applied broadly), the learned filter is the better scorer: on the same broad pools, \texttt{filter} $\times$ \texttt{gen} beats the frequency prior \textbf{0.413 vs 0.374} at \texttt{top\_k}=300 (ranking-signal arm of the selection-breadth ablation above, n=245 subset; \texttt{results/\allowbreak selection\_\allowbreak ablation\_\allowbreak ranksignal.json}), consistent with Proposition 1's listwise-reranker gain. \emph{Falsifiable prediction:} reweighting the labeled loss by $1/\hat{e}(r)$ (a SAR correction) should shrink the prior's edge --- an \textbf{open test, not a promised fix}. \emph{Guardrail:} the propensity model $e(r) \propto \pi(r)$ is an \textbf{unmeasured modeling assumption}, flagged as such --- this is an explanatory model plus a refutable prediction, not proof that learning cannot win, and it is consistent with the deployed pipeline's prior-independence noted above.

\paragraph{Proposition 3 --- Paradigm limit ($\to$ \texttt{coverage\allowbreak\_bank}).}
Single-step rule-based recall $\leq$ single-step \texttt{coverage\allowbreak\_bank} $<1$, because a non-vanishing fraction of references are multi-generation (e.g.\ oxidation $\to$ conjugation) and unreachable by any single rule application --- an irreducible residual that no ranking improvement can recover.

\emph{Witnesses:} depth-2 rule application lifts the ceiling by only \textbf{+0.012} at \textbf{8.5$\times$} candidate cost (Table~\ref{tab:levers}, row 2; \texttt{results/\allowbreak benchmark\_\allowbreak report\_\allowbreak depth2.json}); on the external GLORYx set the uncapped single-step ceiling is \textbf{0.633} (reported in the external-validation section), whose references include single-step-unreachable multi-generation metabolites.

\emph{Composition of the gap:} typing each uncovered transformation by its radius-0 reaction type (element-aware changed-bond multiset) against the \textbf{4,417} types the bank already contains (\texttt{results/\allowbreak coverage\_\allowbreak gap\_\allowbreak types.json}; the pipeline reproduces the ceiling exactly, coverage 0.7355) shows the single-step gap is a mix skewed to novel chemistry --- of \textbf{687} uncovered test transformations, \textbf{41\% (280)} are of a reaction type the bank \emph{already} has (reachable in principle by a more general template, an in-bank ceiling of at most $(1910+280)/2597 = \textbf{0.843}$), while \textbf{53\% (366)} are \textbf{novel types absent from the bank}, addressable only by new reaction corpora (both together would bound the ceiling at 0.984). So template generalization is a real but \textbf{bounded} lever ($\approx$+0.11 ceiling headroom at best); the majority of the residual gap is genuinely out-of-bank chemistry, not a template-specificity artifact. \emph{Guardrail:} the bound is \textbf{single-step-conditional} --- it bounds the single-step paradigm, not the problem; multi-step and out-of-bank chemistry remain open coverage levers.

\subsection{Causal interventions}

\paragraph{Intervention: a factorized generator tests Propositions 2--3.}
This is a controlled intervention turning the observational decomposition above into a causal test: we rebuilt Stage-1 as a dense factorized generator, $P(\text{type}\mid s)\cdot P(\text{site}\mid \text{type},s)$, over a coarse $\sim$radius-0 type vocabulary.

\textbf{Selection is fixable:} the learned type head beats the frequency prior $\sim$3$\times$ at $k=5$ (recall@5 0.086 vs 0.030), robust across two vocabulary granularities (371 and 701 types), converging to the prior by $k=10$ (0.298 vs 0.309); site localisation holds (site hit@3 0.78, cf.\ Gate C 0.81) --- val set, \texttt{results/\allowbreak factorized\_\allowbreak val.json} --- so Proposition 2's PU degeneracy is a target-support artifact, not fundamental, confirming its falsifiable $1/\hat{e}$-style prediction.

\textbf{But coverage binds:} as a replacement generator the pipeline reaches only recall@15 0.256 on the full clean test --- a net regression against deployed GRAIL (budget-matched recall@15 0.365, precision 0.109) --- because type-selection caps oracle reachability at $\sim$47\%, unchanged when the vocabulary widened 371$\to$701 (\texttt{results/\allowbreak factorized\_\allowbreak eval.json}), confirming Proposition 3 causally.

\textbf{Ranking margin:} kept at broad coverage and used only to re-rank, \texttt{filter}$\times$\texttt{gen}$\times$\texttt{type}$\times$\texttt{site} lifts recall@15 from 0.388 to 0.404 over \texttt{filter}$\times$\texttt{gen} on the identical broad pool --- a paired +0.0165 (95\% CI [+0.006, +0.027], paired bootstrap over the full clean test, n=1170); the factorized signal alone is n.s.\ (+0.001, CI [$-$0.010, +0.013]) (\texttt{results/\allowbreak hybrid\_\allowbreak rerank\_\allowbreak full1170.json}).

\textbf{Joint vs.\ bolt-on ranking:} the bolt-on re-ranker multiplies \emph{independently} MLE-trained type/site probabilities into the rank; fine-tuning those two heads instead with a listwise ranking loss against the \textbf{fixed} generator and filter (frozen $\log(\text{filter})+\log(\text{gen})$ context, gradients only through the heads) sharpens the signal --- on a matched \texttt{top\_k}=100 pool the joint-trained re-ranker beats the bolt-on by a paired \textbf{+0.0089 (95\% CI [+0.003, +0.015], n=1170)} and the \texttt{filter}$\times$\texttt{gen} baseline by \textbf{+0.0091 [+0.003, +0.016]}, where on that pool the independently-trained bolt-on adds nothing (+0.0002, n.s.). So \emph{rank-aware} training, not the type/site heads per se, recovers this margin (\texttt{results/\allowbreak joint\_\allowbreak rerank.json}).

\emph{Synthesis:} selection is addressable, ranking yields a small significant margin, but coverage --- the source-limited rule bank --- remains binding; no Stage-1 redesign matches SyGMa without a broader bank.

\paragraph{Set-generation (GFlowNet): a released negative result.}
Beyond re-ranking, we also formulated Stage-1 as a \textbf{set-terminal GFlowNet} (\texttt{model/\allowbreak set\_\allowbreak gflownet.py}): the terminal object is a forest/set of metabolites per substrate, trained by Trajectory Balance against a PU set-coverage reward with an analytic \texttt{1/\#leaves} backward --- a genuinely novel formulation for this domain, released for reuse. On our benchmark it does \textbf{not} beat the simpler baselines: across 3 seeds (val, n=40) its single-set recall@15 is \textbf{0.292$\pm$0.032}, statistically indistinguishable from beam (0.311$\pm$0.039) and below the Stage-2a reranker (0.411$\pm$0.044). It does produce diverse sets (pairwise-Tanimoto 0.214$\pm$0.003, $\sim$39 unique scaffolds), but a clean \emph{diversity-advantage} claim is \textbf{not} supported here --- the union-coverage comparison against the baselines is dominated by an under-production artefact in the baseline arms (its variance spans 0.0--0.58 across seeds), so we do not report that metric. We include this as a negative result because a benchmark is well-served by an honest report that a fashionable architecture does not win on it (\texttt{results/gflownet\_seed\{0,1,2\}\_overnight.json}; a fair multi-seed diversity comparison at scale is future work, discussed in the future-work section).

\bigskip
\emph{Source: \texttt{results/\allowbreak prior\_\allowbreak vs\_\allowbreak learned.json}, \texttt{results/\allowbreak benchmark\_\allowbreak report\_\allowbreak depth2.json}, \texttt{results/\allowbreak benchmark\_\allowbreak report\_\allowbreak gap.json}, \texttt{docs/\allowbreak benchmark/\allowbreak stage2\_\allowbreak ranker\_\allowbreak evidence.md}, \texttt{results/\allowbreak selection\_\allowbreak ablation.json}, \texttt{results/\allowbreak selection\_\allowbreak ablation\_\allowbreak ranksignal.json}.}
